\hypertarget{ordinary-differential-equations}{%
\subsection{Ordinary Differential
Equations}\label{sec:odes}}

Ordinary differential equations offer a powerful framework for
formulating hysteretic models.
For example, one-dimensional perfect plasticity can be formulated as the following ODE:
% \[
% \begin{aligned}
% \dot{\sigma} &=E\left(1-\frac{1+\operatorname{sgn}(\sigma \dot{\varepsilon})}{2}\left\{U\left(\sigma_{y}^{+}-\sigma\right) U(\sigma)
% +U\left(-\sigma_{y}^{-}+\sigma\right) U(-\sigma)\right\}\right) \dot{\varepsilon}
% \end{aligned}
% \]
\begin{equation}
\begin{aligned}
\dot{\bar{\sigma}} &= \phi(\bar{\sigma}, \dot{\varepsilon}) \, \dot{\varepsilon}\\
\phi(\bar{\sigma}, \dot{\varepsilon})&=\left\{\begin{array}{lll}
E & \text { if } | \bar{\sigma}|<1 & \text { or } \, \, \, \, \, \, \, \bar{\sigma} \operatorname{sgn}(\dot{\varepsilon}) \leq 0 \\
bE & \text { if }|\bar{\sigma}|\geq 1 & \text { and } \, \, \, \, \bar{\sigma} \operatorname{sgn}(\dot{\varepsilon})>0
\end{array}\right.
\end{aligned}
\label{eq:ode_general}
\end{equation}
where $\bar{\sigma}=\sigma/\sigma_y$, $\sigma_y$ is the yield stress, $E$ is the elastic modulus, $b$ is the strain hardening parameter, $\operatorname{sgn}$ is the signum function{\footnote{The signum function of a real number $x$ is a piecewise function which is defined as follows:
\[
\operatorname{sgn}(x) \triangleq \left\{\begin{array}{cl}
-1 & \text { if } x<0, \\
0 & \text { if } x=0, \\
1 & \text{ if } x>0.
\end{array}
\right.
\]}}. 
An important family of hysteretic models (with various forms of the function \(\phi\)) that generalizes this formulation are known as Bouc-Wen models and are generally given by a relation of the following form:
\begin{equation}
\begin{aligned}
\sigma(\varepsilon) &= b E \varepsilon+(1-b) E z(\varepsilon),
\\
\dot{z}&=\frac{\partial z}{\partial \varepsilon} \frac{\partial \varepsilon}{\partial t} = \phi(z,\varepsilon,\dot \varepsilon) \, \dot{\varepsilon}.
\end{aligned}
\label{eq:bouc-family}
\end{equation}

\begin{description}
\tightlist
\item[Bouc's Hysteresis (1967)]
The original model presented by \cite{bouc1967forced} can be given in the form of \cref{eq:bouc-family} with \(\phi\) given by: 
\begin{equation}
\phi(z,\dot{\varepsilon})=A-\alpha \operatorname{sgn}(\dot{\varepsilon}) \, z,
%\dot{z}=(A-\alpha \operatorname{sgn}(\dot{\varepsilon}) z) \dot{\varepsilon} \\
\label{eq:bouc-family1}
\end{equation}
where $A$ scales the hysteretic loop and  $\alpha$ controls the 
shape of the hysteresis loop.
Bouc's model yields a separable initial-value problem for \(\alpha \ne 0\) with solutions:
\begin{equation}
  z(\varepsilon) = \begin{cases}-\frac{1}{\alpha}\left(C_1 e^{-\alpha \varepsilon}-A\right), & \text { if } \dot{\varepsilon}>0 \\ \frac{1}{\alpha}\left(C_2 e^{\alpha \varepsilon}-A\right), & \text { if } \dot{\varepsilon}<0\end{cases}
\label{eq:bouc-family2}
\end{equation}
where \(C_1\) and \(C_2\) must be calculated piecewise by using the
final state of the previous branch as the initial condition.

\item[Wen's Hysteresis (1976)]
Wen \cite{wen1976method} extended the original model by Bouc to control the sharpness of the hysteresis in transition from elastic to inelastic region:
\begin{equation}
%\dot{z}=\left(A-(\alpha \operatorname{sgn}(\dot{\varepsilon} z)-\beta) \, |z|^{n}\right) \dot{\varepsilon} \\
\phi(z,\dot{\varepsilon}) = A-\left(\alpha \operatorname{sgn}(z \dot{\varepsilon})-\beta\right)|z|^{n},
\label{eq:bouc-family3}
\end{equation}
where $n$ is a parameter that controls the sharpness.

\item[Baber-Wen Degradation (1981)]
Baber and Wen (1981) \cite{baber1981random} further extended the model to 
include strength and stiffness deterioration:
\begin{equation}
\begin{aligned}
\phi(z,\dot{\varepsilon})&=\frac{1}{\eta}\left(A- \nu |z|^{n} \left(\alpha \operatorname{sgn}(\dot{\varepsilon} z)+\beta\right) \right),
\\
A &=A_{o}-\delta_{A} \psi,   \\
\nu &=1+\delta_{\nu} \psi,   \\
\eta &=1+\delta_{\eta} \psi, \\
\dot{\psi} & =(1-b) E \dot{\varepsilon} z, 
\end{aligned}
\label{eq:bouc-family4}
\end{equation}
where $A_{o}$, $\delta_{A}$, $\delta_{\nu}$, and $\delta_{\eta}$ are degradation parameters and $\psi$ is the absorbed hysteretic energy, defined as follows:
\begin{equation}
\psi = (1-b) \omega^2 \int_0^t \dot{\varepsilon}(\tau) z(\tau) \mathrm{d} \tau,
\label{eq:bouc-family5}
\end{equation}
where $\rho$ is the material density and $\omega^2:={E}/{\rho}$ is the squared pseudo-natural frequency of the nonlinear system.
% \item[Noori]
% \[
% \phi =\frac{h(z)}{\eta} \bigl(A-\nu\left(\beta \operatorname{sign}(\dot{u})|z|^{n-1} z+\gamma |z|^n\right)\bigr)
% \]
% \item[Wang and Wen (1998)] suggested the following expression to account for the asymmetric peak restoring force:
% \[
% \phi(z, \dot{u}) = A - |z|^n \left(\gamma+\beta \operatorname{sgn}(z \dot{u})+\varphi(\operatorname{sgn}(\dot{u})+\operatorname{sgn}(z))\right)
% \]
\end{description}

\hypertarget{implementation}{%
\subsection{Implementations}\label{implementations}}
When such models are employed within a finite element simulation,
the governing ODE for the internal variables must be integrated numerically. For example, applying the implicit \emph{Backward-Euler} discretization to the variables \(\dot{\psi}\) and \(\dot z\) from Bouc-Wen-Baber (1981) \cite{baber1981random} furnishes the
following implicit problem for \(z_{n+1}\) and \(\psi_{n+1}\):

\begin{quote}
\[
\begin{aligned}
z_{n+1}&=z_{n}+\Delta t \frac{A_{n+1}-\left|z_{n+1}\right|^{n}\left\{\gamma+\beta \operatorname{sgn}\left(\frac{\epsilon_{n+1}-\epsilon_{n}}{\Delta t} z_{n+1}\right)\right\} \nu_{n+1}}{\eta_{n+1}} \, \frac{\epsilon_{n+1}-\epsilon_{n}}{\Delta t}, \\
\psi_{n+1}
&=\psi_{n}+\Delta t(1-\alpha) k_{o} \frac{\epsilon_{n+1}-\epsilon_{n}}{\Delta t} z_{n+1}, \\
A_{n+1} &=A_{o}-\delta_{A} \, \psi_{n+1}, \\
\nu_{n+1} &=1+\delta_{\nu} \, \psi_{n+1}, \\
\eta_{n+1} &=1+\delta_{\eta} \, \psi_{n+1},
\end{aligned}
\]
\end{quote}
which is generally solved with Newton-Raphson iterations.
